\begin{abstractd}
Este Proyecto Integrador de Aprendizaje presenta la resolución completa de las siete integrales incluidas en la consigna. Aunque la evaluación presencial solicita dos ejercicios al azar, se desarrollan todos como guía de preparación. Se emplean sustitución algebraica, expansión de productos notables y completación de cuadrados; cada resultado se verifica mediante derivación.
\end{abstractd}

\templateIndex

\section{Introducción}
La integración indefinida consiste en determinar una familia de antiderivadas. Si $F'(x)=f(x)$, entonces
\[
 \int f(x)\,dx=F(x)+C,
\]
donde $C$ es una constante. La selección del método depende de la estructura del integrando: una expresión compuesta acompañada por su derivada sugiere sustitución; un cuadrado algebraico puede expandirse; y un denominador cuadrático puede simplificarse completando el cuadrado \citep{stewart2016calculus,thomas2018calculus}.

\section{Resumen de métodos}
Los criterios utilizados son los siguientes:
\begin{itemize}
  \item \textbf{Sustitución:} usar $u=g(x)$, $du=g'(x)\,dx$ cuando un cambio de variable convierte el integrando en una potencia elemental.
  \item \textbf{Simplificación algebraica:} expandir o reescribir antes de integrar término a término.
  \item \textbf{Cuadrado perfecto:} transformar $x^2+8x-7$ en $(x+4)^2-23$ para reconocer una derivada directa.
  \item \textbf{Verificación:} derivar la antiderivada obtenida y recuperar exactamente el integrando original.
\end{itemize}

\section{Resolución de ejercicios}

\begin{activitybox}{Ejercicio 1: $\displaystyle \int x\sqrt{4x+1}\,dx$}
\textbf{Técnica:} sustitución lineal. Sea
\[
 u=4x+1,\qquad du=4\,dx,\qquad dx=\frac{du}{4},\qquad x=\frac{u-1}{4}.
\]
Entonces
\begin{align*}
 \int x\sqrt{4x+1}\,dx
 &=\int \frac{u-1}{4}u^{1/2}\frac{du}{4}\\
 &=\frac1{16}\int\left(u^{3/2}-u^{1/2}\right)du\\
 &=\frac1{16}\left(\frac{u^{5/2}}{5/2}-\frac{u^{3/2}}{3/2}\right)+C\\
 &=\frac{u^{5/2}}{40}-\frac{u^{3/2}}{24}+C.
\end{align*}
Al regresar a $x$ y factorizar la potencia común $(4x+1)^{3/2}$,
\begin{align*}
 \frac{(4x+1)^{5/2}}{40}-\frac{(4x+1)^{3/2}}{24}+C
 &=(4x+1)^{3/2}\left(\frac{4x+1}{40}-\frac1{24}\right)+C\\
 &=(4x+1)^{3/2}\left(\frac{3(4x+1)-5}{120}\right)+C\\
 &=\frac{(6x-1)(4x+1)^{3/2}}{60}+C.
\end{align*}
Por tanto,
\[
 \boxed{\int x\sqrt{4x+1}\,dx
 =\frac{(6x-1)(4x+1)^{3/2}}{60}+C}.
\]
\textbf{Verificación:}
\[
 \frac{d}{dx}\left[\frac{(6x-1)(4x+1)^{3/2}}{60}\right]
 =\frac{(4x+1)^{1/2}}{60}\bigl[6(4x+1)+6(6x-1)\bigr]
 =x\sqrt{4x+1}.
\]
\end{activitybox}

\begin{activitybox}{Ejercicio 2: $\displaystyle \int(x+1)\sqrt{2-x}\,dx$}
\textbf{Técnica:} sustitución lineal. Se toma
\[
 u=2-x,\qquad du=-dx,\qquad x=2-u,
\]
de modo que $x+1=3-u$. Por tanto,
\begin{align*}
 \int(x+1)\sqrt{2-x}\,dx
 &=-\int(3-u)u^{1/2}\,du\\
 &=\int\left(u^{3/2}-3u^{1/2}\right)du\\
 &=\frac{u^{5/2}}{5/2}-3\frac{u^{3/2}}{3/2}+C\\
 &=\frac25u^{5/2}-2u^{3/2}+C.
\end{align*}
Sustituyendo $u=2-x$ y factorizando,
\begin{align*}
 \frac25(2-x)^{5/2}-2(2-x)^{3/2}+C
 &=(2-x)^{3/2}\left[\frac25(2-x)-2\right]+C\\
 &=-\frac25(x+3)(2-x)^{3/2}+C.
\end{align*}
Así,
\[
 \boxed{\int(x+1)\sqrt{2-x}\,dx
 =-\frac{2(x+3)(2-x)^{3/2}}{5}+C}.
\]
\textbf{Verificación:}
\begin{align*}
 \frac{d}{dx}\left[-\frac25(x+3)(2-x)^{3/2}\right]
 &=-\frac25(2-x)^{3/2}+\frac35(x+3)(2-x)^{1/2}\\
 &=(2-x)^{1/2}\left[-\frac25(2-x)+\frac35(x+3)\right]\\
 &=(x+1)\sqrt{2-x}.
\end{align*}
\end{activitybox}

\begin{activitybox}{Ejercicio 3: $\displaystyle \int\frac{2x+1}{\sqrt{x+4}}\,dx$}
\textbf{Técnica:} sustitución lineal. Sea
\[
 u=x+4,\qquad du=dx,\qquad x=u-4,
\]
por lo que $2x+1=2u-7$. Así,
\begin{align*}
 \int\frac{2x+1}{\sqrt{x+4}}\,dx
 &=\int\frac{2u-7}{u^{1/2}}\,du\\
 &=\int\left(2u^{1/2}-7u^{-1/2}\right)du\\
 &=2\frac{u^{3/2}}{3/2}-7\frac{u^{1/2}}{1/2}+C\\
 &=\frac43u^{3/2}-14u^{1/2}+C.
\end{align*}
Al sustituir $u=x+4$ y extraer el factor $\sqrt{x+4}$,
\begin{align*}
 \frac43(x+4)^{3/2}-14(x+4)^{1/2}+C
 &=\sqrt{x+4}\left[\frac43(x+4)-14\right]+C\\
 &=\frac{2(2x-13)\sqrt{x+4}}3+C.
\end{align*}
Por consiguiente,
\[
 \boxed{\int\frac{2x+1}{\sqrt{x+4}}\,dx
 =\frac{2(2x-13)\sqrt{x+4}}{3}+C}.
\]
\textbf{Verificación:}
\begin{align*}
 \frac{d}{dx}\left[\frac23(2x-13)(x+4)^{1/2}\right]
 &=\frac43\sqrt{x+4}+\frac{2x-13}{3\sqrt{x+4}}\\
 &=\frac{4(x+4)+2x-13}{3\sqrt{x+4}}
 =\frac{2x+1}{\sqrt{x+4}}.
\end{align*}
\end{activitybox}

\begin{activitybox}{Ejercicio 4: $\displaystyle \int t\sqrt[3]{t+10}\,dt$}
\textbf{Técnica:} sustitución lineal. Se define
\[
 u=t+10,\qquad du=dt,\qquad t=u-10.
\]
Entonces
\begin{align*}
 \int t\sqrt[3]{t+10}\,dt
 &=\int(u-10)u^{1/3}\,du\\
 &=\int\left(u^{4/3}-10u^{1/3}\right)du\\
 &=\frac{u^{7/3}}{7/3}-10\frac{u^{4/3}}{4/3}+C\\
 &=\frac37u^{7/3}-\frac{15}{2}u^{4/3}+C.
\end{align*}
Al volver a la variable original y factorizar $(t+10)^{4/3}$,
\begin{align*}
 \frac37(t+10)^{7/3}-\frac{15}{2}(t+10)^{4/3}+C
 &=(t+10)^{4/3}\left[\frac37(t+10)-\frac{15}{2}\right]+C\\
 &=\frac{3(2t-15)(t+10)^{4/3}}{14}+C.
\end{align*}
En consecuencia,
\[
 \boxed{\int t\sqrt[3]{t+10}\,dt
 =\frac{3(2t-15)(t+10)^{4/3}}{14}+C}.
\]
\textbf{Verificación:}
\begin{align*}
 \frac{d}{dt}\left[\frac{3}{14}(2t-15)(t+10)^{4/3}\right]
 &=\frac37(t+10)^{4/3}+\frac27(2t-15)(t+10)^{1/3}\\
 &=(t+10)^{1/3}\frac{3(t+10)+2(2t-15)}7\\
 &=t\sqrt[3]{t+10}.
\end{align*}
\end{activitybox}

\begin{activitybox}{Ejercicio 5: $\displaystyle \int\left(x+\frac1x\right)^2dx$}
\textbf{Técnica:} expansión algebraica e integración término a término. Para $x\neq0$,
\begin{align*}
 \left(x+\frac1x\right)^2
 &=x^2+2x\left(\frac1x\right)+\left(\frac1x\right)^2\\
 &=x^2+2+x^{-2}.
\end{align*}
Por la regla de la potencia,
\begin{align*}
 \int\left(x+\frac1x\right)^2dx
 &=\int\left(x^2+2+x^{-2}\right)dx\\
 &=\frac{x^{2+1}}{2+1}+2x+\frac{x^{-2+1}}{-2+1}+C\\
 &=\frac{x^3}{3}+2x-x^{-1}+C\\
 &=\frac{x^3}{3}+2x-\frac1x+C.
\end{align*}
Por tanto,
\[
 \boxed{\int\left(x+\frac1x\right)^2dx
 =\frac{x^3}{3}+2x-\frac1x+C},\qquad x\neq0.
\]
\textbf{Verificación:}
\[
 \frac{d}{dx}\left(\frac{x^3}{3}+2x-\frac1x\right)
 =x^2+2+\frac1{x^2}
 =\left(x+\frac1x\right)^2.
\]
\end{activitybox}

\begin{activitybox}{Ejercicio 6: $\displaystyle \int3x^2\sqrt{2x^3-5}\,dx$}
\textbf{Técnica:} sustitución por la función interna. Sea
\[
 u=2x^3-5,\qquad du=6x^2\,dx,
 \qquad 3x^2\,dx=\frac12du.
\]
Así,
\begin{align*}
 \int3x^2\sqrt{2x^3-5}\,dx
 &=\frac12\int u^{1/2}\,du\\
 &=\frac12\left(\frac{u^{1/2+1}}{1/2+1}\right)+C\\
 &=\frac12\left(\frac23u^{3/2}\right)+C\\
 &=\frac13u^{3/2}+C.
\end{align*}
Finalmente se sustituye $u=2x^3-5$; en consecuencia,
\[
 \boxed{\int3x^2\sqrt{2x^3-5}\,dx
 =\frac13(2x^3-5)^{3/2}+C}.
\]
\textbf{Verificación:}
\[
 \frac{d}{dx}\left[\frac13(2x^3-5)^{3/2}\right]
 =\frac13\cdot\frac32(2x^3-5)^{1/2}(6x^2)
 =3x^2\sqrt{2x^3-5}.
\]
\end{activitybox}

\begin{activitybox}{Ejercicio 7: $\displaystyle \int\frac{x+4}{(x^2+8x-7)^2}\,dx$}
\textbf{Técnica:} completación de cuadrados y sustitución. El denominador se reescribe como
\begin{align*}
 x^2+8x-7
 &=x^2+8x+16-16-7\\
 &=(x+4)^2-23.
\end{align*}
Sea
\[
 u=x+4,\qquad du=dx.
\]
Entonces
\[
 \int\frac{x+4}{(x^2+8x-7)^2}\,dx
 =\int\frac{u}{(u^2-23)^2}\,du.
\]
Con una segunda sustitución,
\[
 v=u^2-23,\qquad dv=2u\,du,
 \qquad u\,du=\frac12dv,
\]
se obtiene
\begin{align*}
 \int\frac{u}{(u^2-23)^2}\,du
 &=\frac12\int v^{-2}\,dv\\
 &=\frac12\left(\frac{v^{-1}}{-1}\right)+C\\
 &=-\frac1{2v}+C.
\end{align*}
Al deshacer ambas sustituciones, $v=u^2-23=(x+4)^2-23=x^2+8x-7$. Por lo tanto,
\[
 \boxed{\int\frac{x+4}{(x^2+8x-7)^2}\,dx
 =-\frac1{2(x^2+8x-7)}+C}.
\]
\textbf{Verificación:}
\[
 \frac{d}{dx}\left[-\frac12(x^2+8x-7)^{-1}\right]
 =\frac12(2x+8)(x^2+8x-7)^{-2}
 =\frac{x+4}{(x^2+8x-7)^2}.
\]
\end{activitybox}

\section{Resultados finales}
Las antiderivadas obtenidas se concentran en la Tabla~\ref{tab:resultados-pia}.
\begin{table}[H]
\centering
\caption{Resumen de resultados del PIA.}
\label{tab:resultados-pia}
\small
\begin{tabular}{cl}
\toprule
Ejercicio & Antiderivada \\
\midrule
1 & $\dfrac{(6x-1)(4x+1)^{3/2}}{60}+C$ \\[2mm]
2 & $-\dfrac{2(x+3)(2-x)^{3/2}}5+C$ \\[2mm]
3 & $\dfrac{2(2x-13)\sqrt{x+4}}3+C$ \\[2mm]
4 & $\dfrac{3(2t-15)(t+10)^{4/3}}{14}+C$ \\[2mm]
5 & $\dfrac{x^3}{3}+2x-\dfrac1x+C$ \\[2mm]
6 & $\dfrac13(2x^3-5)^{3/2}+C$ \\[2mm]
7 & $-\dfrac1{2(x^2+8x-7)}+C$ \\
\bottomrule
\end{tabular}
\end{table}

\section{Conclusiones}
La estructura algebraica de cada integral permitió seleccionar un procedimiento breve y verificable. Los ejercicios 1 a 4 se redujeron a potencias mediante sustituciones lineales; el ejercicio 5 requirió expandir un binomio; el ejercicio 6 exhibió directamente la derivada de su función interna; y el ejercicio 7 se simplificó al completar el cuadrado. La derivación de cada resultado recuperó el integrando correspondiente, lo que confirma las siete soluciones. La práctica completa también proporciona preparación para resolver correctamente cualquier pareja seleccionada al azar durante la evaluación presencial.

\nocite{openstax2016calculus1,openstax2016calculus2}
\bibliography{UANL/ingeniero-agronomo/calculo-integral/calculo-integral}
